% Linux kernel validation of the organization column (needs booktabs; cites paper-snippet.bib).
% Numbers from linux-validation-20261107/RESULTS.md.
\paragraph{External check on the Linux kernel.}
We compare the organization shares of the Linux kernel (pinned at v7.3-rc7) with the employer
statistics that LWN.net publishes for each release~\cite{corbet2026dev73,corbet2025dev618} and
with the cumulative 2007--2019 table of the Linux Foundation~\cite{stewart2020history}. These
sources credit each commit to the employer of its author through gitdm, a map of address domains
that the editors complete by hand~\cite{corbet2007who}. No public source splits the code alive in
the tree by employer, so we match periods by the committer date of the blamed commit
(Table~\ref{tab:linux-validation}). Over the employers that each source lists, the Spearman
correlation is $\rho=0.84$ ($n=16$, permutation $p<0.001$) between our surviving commits and
LWN's changesets for 7.3, $\rho=0.73$ ($n=17$, $p=0.001$) between our alive tokens and LWN's lines
changed for 7.3, and $\rho=0.65$ ($n=17$, $p=0.005$) between our surviving commits and the Linux
Foundation's commits for 2007--2019. The mean absolute differences over each source's top~10 are
1.2, 2.2 and 0.9 percentage points. The share without an organization is close: 16.9\% of our
7.3 tokens against 18.6\% of LWN's lines, where LWN counts its (Unknown), (None) and (Consultant)
rows. The remaining gaps have known causes. First, lines changed count deletions, which leave no
alive tokens. Second, AMD's generated register headers are 11.4\% of the alive code tokens. Third,
the address rule: the project domain \texttt{kernel.org}, which the dataset maps to the Linux
Foundation, carries 8.1\% of the 7.3 tokens. LWN credits these maintainers to their employers by
hand, which likely explains our low shares for Google and Meta.

% For a two-column layout, use table* or keep \footnotesize.
\begin{table}[t]
\centering
\footnotesize
\setlength{\tabcolsep}{3.5pt}
\caption{Linux kernel: shares (\%) by organization in the dataset (ours) and in independent
employer statistics. Ours counts the code tokens alive at v7.3-rc7, or the commits that still own
an alive token, whose committer date falls in the source period. The sources count all
changesets, lines changed or commits of the period. n.l.: not in the source's top~20.}
\label{tab:linux-validation}
\begin{tabular}{@{}lrrrrrrr@{}}
\toprule
 & \multicolumn{1}{c}{All} & \multicolumn{4}{c}{7.3 cycle} & \multicolumn{2}{c}{2007--2019} \\
\cmidrule(lr){2-2}\cmidrule(lr){3-6}\cmidrule(l){7-8}
 & Ours & Ours & LWN & Ours & LWN & Ours & LF \\
Organization & tokens & tokens & lines & commits & changesets & commits & commits \\
\midrule
AMD              & 16.3 & 29.4 & 22.3 & 5.7  & 6.3  & 2.6 & 2.3  \\
Intel            & 7.0  & 3.3  & 4.4  & 5.4  & 6.3  & 9.0 & 10.0 \\
Red Hat          & 3.8  & 1.4  & 3.8  & 2.4  & 4.5  & 6.9 & 8.9  \\
Linux Foundation & 3.6  & 8.1  & n.l. & 12.6 & n.l. & 2.7 & 0.8  \\
Qualcomm         & 3.0  & 6.4  & 6.1  & 3.6  & 5.9  & 1.5 & n.l. \\
Realtek          & 2.5  & 2.2  & 1.6  & 0.9  & n.l. & 0.2 & n.l. \\
Linaro           & 2.1  & 0.1  & n.l. & 0.3  & n.l. & 2.3 & 3.2  \\
NVIDIA           & 1.6  & 1.3  & 1.7  & 2.1  & 2.1  & 0.8 & n.l. \\
IBM              & 1.5  & 0.9  & 2.0  & 2.2  & 2.5  & 2.8 & 3.8  \\
Broadcom         & 1.4  & 0.2  & n.l. & 0.3  & n.l. & 1.0 & 1.2  \\
Google           & 1.3  & 2.0  & 5.6  & 3.8  & 6.1  & 1.6 & 2.8  \\
Meta             & 0.7  & 0.1  & 2.7  & 0.4  & 2.8  & 0.4 & n.l. \\
\midrule
No organization  & 17.4 & 16.9 & 18.6 & 34.7 & 27.1 & 28.5 & 19.0 \\
\midrule
Spearman $\rho$ (listed) & & \multicolumn{2}{c}{0.73 ($n=17$)} & \multicolumn{2}{c}{0.84 ($n=16$)} & \multicolumn{2}{c}{0.65 ($n=17$)} \\
\bottomrule
\end{tabular}
\end{table}
